% =====================================================================
%  APÊNDICE A — GUIA DE INSTALAÇÃO E EXECUÇÃO
% =====================================================================

\chapter{GUIA DE INSTALAÇÃO E EXECUÇÃO}
\label{ap:instalacao}

Este apêndice apresenta os pré-requisitos e os passos para instalação e
execução do \argus{} em modo de desenvolvimento.

\section{Pré-requisitos}
\label{ap:instalacao-pre-requisitos}

\begin{itemize}
    \item Go (v1.20 ou superior) --- \url{https://go.dev/};
    \item Node.js e npm;
    \item Wails CLI --- \code{go install github.com/wailsapp/wails/v2/cmd/wails@latest};
    \item Android Studio (apenas para a compilação móvel via Capacitor).
\end{itemize}

\section{Instalação das dependências}
\label{ap:instalacao-dependencias}

\begin{enumerate}
    \item Clonar o repositório e entrar na pasta do projeto;
    \item Instalar as dependências do \textit{backend} Go:
        \begin{verbatim}
        go mod tidy
        \end{verbatim}
    \item Instalar as dependências do \textit{frontend}:
        \begin{verbatim}
        cd frontend
        npm install
        cd ..
        \end{verbatim}
\end{enumerate}

\section{Execução em modo de desenvolvimento}
\label{ap:instalacao-execucao}

Com as dependências instaladas, a aplicação pode ser executada em modo de
desenvolvimento (com \textit{hot reload}) por meio do comando:

\begin{verbatim}
wails dev
\end{verbatim}

\complementar{Descrever aqui, se desejado, o fluxo de geração do instalador de
produção para Windows/Linux (\code{wails build}) e o processo de compilação do
aplicativo Android com Capacitor (\code{npx cap sync android}), incluindo
capturas de tela do passo a passo.}
